\section{Estadistica descriptiva}
Antes de profundizar en el tema es importante tener una forma de dimensionar los datos con
los que estamos trabajando mediante medidas de tendencia central y dispersión. A continuación, se
presentan los momentos estadísticos que permiten entender la distribución global de las muestras.

\begin{table}[H]
\centering
\caption{Estadísticas Descriptivas por sujeto control (Parte 1/2)}
\label{tab:estadisticas1}
\begin{tabular}{lSSSSS}
\toprule
\textbf{Métrica} & \textbf{EMV} & \textbf{GH2} & \textbf{GUR} & \textbf{JAL} & \textbf{JAN} \\
\midrule
count & 8.06e6 & 8.06e6 & 8.06e6 & 8.06e6 & 8.06e6 \\
mean  & -0.192 & 1.23  & -0.982 & -1.11 & -0.731 \\
std   & 5.99  & 3.24  & 5.95  & 3.81 & 4.50  \\
min   & -60.5 & -43.7 & -98.2 & -40.0 & -55.7 \\
25\%  & -3.10 & -0.400 & -2.30 & -2.50 & -1.70 \\
50\%  & -0.400 & 1.10 & -0.900 & -1.00 & 0.00 \\
75\%  & 3.00 & 2.70 & 0.500 & 0.300 & 1.50 \\
max   & 209 & 60.4 & 61.2 & 92.8 & 94.2 \\
\bottomrule
\end{tabular}
\end{table}

\begin{table}[h]
\centering
\caption{Estadísticas Descriptivas por sujeto control (Parte 2/2)}
\label{tab:estadisticas2}
\begin{tabular}{lSSSSS}
\toprule
\textbf{Métrica} & \textbf{MGN} & \textbf{MJN} & \textbf{MMA} & \textbf{RAN} & \textbf{VCN} \\
\midrule
count & 8.06e6 & 8.06e6 & 8.06e6 & 8.06e6 & 8.06e6 \\
mean  & 0.279  & -2.71 & -2.63 & -1.70 & -3.56 \\
std   & 3.65   & 4.61  & 1.84  & 4.12  & 2.64  \\
min   & -33.7  & -48.6 & -21.6 & -49.2 & -48.0 \\
25\%  & -1.30  & -5.10 & -3.50 & -3.50 & -4.80 \\
50\%  & 0.500  & -2.80 & -2.60 & -1.70 & -3.40 \\
75\%  & 2.10   & -0.500 & -1.70 & 0.00  & -2.20 \\
max   & 46.2   & 98.2  & 33.5  & 60.3  & 81.1  \\
\bottomrule
\end{tabular}
\end{table}

\begin{table}[h]
\centering
\caption{Estadísticas Descriptivas por sujeto control}
\label{tab:estadisticas3}
\begin{tabular}{lSSSSSSSS}
\toprule
\textbf{Métrica} & \textbf{AEF} & \textbf{CLM} & \textbf{FGV} & \textbf{JGM} & \textbf{LIV} & \textbf{PCM} & \textbf{RLM} & \textbf{RRM} \\
\midrule
count & 8.06e6 & 8.06e6 & 8.06e6 & 8.06e6 & 8.06e6 & 8.06e6 & 8.06e6 & 8.06e6 \\
mean  & -1.49 & -2.56 & -1.89 & -0.340 & -9.83 & 1.66 & -1.67 & -13.3 \\
std   & 4.02  & 4.07  & 2.82  & 5.50   & 4.45  & 4.69 & 6.26  & 3.94  \\
min   & -62.7 & -58.5 & -36.3 & -45.9  & -51.2 & -98.5 & -132  & -65.5 \\
25\%  & -2.50 & -3.80 & -3.40 & -2.50  & -11.5 & -0.100 & -2.80 & -14.6 \\
50\%  & -1.40 & -2.40 & -2.20 & -0.500 & -10.3 & 1.30  & -1.50 & -13.8 \\
75\%  & -0.400 & -1.00 & -0.700 & 1.50   & -8.60 & 2.90  & -0.100 & -12.9 \\
max   & 336   & 60.5  & 53.4  & 96.4   & 38.9  & 108   & 410   & 53.4  \\
\bottomrule
\end{tabular}
\end{table}

En ambas tablas (Cuadro 1 y Cuadro 2) se presentan los momentos estadísticos listados a
continuación:
\begin{itemize}
  \item count: refiriéndose a la cantidad de datos con las que cuenta la muestra. 
  \item mean: refiriéndose a la media de los datos.
   \item std: refiriéndose a la desviación estándar.
   \item min: refiriéndose al valor mínimo que presenta la muestra.
   \item 25 \%: refiriéndose al primer cuartil.
   \item 50 \%: refiriéndose al segundo cuartil.
   \item 75 \%: refiriéndose al tercer cuartil.
   \item max: refiriéndose al valor máximo que presenta la muestra.
\end{itemize}


